\section{Introduction}
\label{sec:intro}

Foundation models pre-trained at scale on medical images or image--text pairs have become a common
starting point for medical image analysis. Self-supervised encoders for
radiology~\cite{perezgarcia2024scalable}, computational
pathology~\cite{chen2024generalpurpose,xu2024wholeslide,vorontsov2024model} and
ophthalmology~\cite{zhou2023model}, most of them trained by the DINOv2 recipe~\cite{oquab2024dinov2}, are
reused frozen across dozens of tasks. Dual-encoder vision--language models (VLMs) built on the CLIP
recipe~\cite{radford2021learning}, such as BiomedCLIP~\cite{zhang2023biomedclip} and
CONCH~\cite{lu2024visuallanguage}, support zero-shot classification and retrieval. Multimodal large
language models (MLLMs) such as LLaVA-Med~\cite{li2023llavamed} and MedGemma~\cite{sellergren2025medgemma}
answer questions about images, and MAIRA-2~\cite{bannur2024maira2} drafts reports
(Table~\ref{tab:models}, Fig.~\ref{fig:overview}). Reuse across institutions, modalities and tasks is what
makes them foundation models~\cite{bommasani2021opportunities}.

Because the same representation is reused so widely, what it encodes matters on its own. Self-reported race
is decodable from medical images~\cite{banerjee2021ai} and from a frozen foundation-model
encoder~\cite{zheng2024demographic}, and is therefore available to any downstream model that reads those
features. A chest-radiograph vision--language system can answer accurately without the
image~\cite{lotfinia2026vision}. A clinically named concept that is present in a representation need not
be used~\cite{basu2026interpretability}. Regulators now ask developers to declare the sources of bias and
the dependencies of their systems~\cite{fda2025aiguidance,euaiact2024}. Input-pixel attribution, the
dominant form of explanation in medical imaging, does not describe what the features contain, and it fails
faithfulness tests in medical-image protocols~\cite{arun2021trustworthiness,sadanandan2026attention}.
Feature-space interpretability addresses the representations themselves.

The methods are not new. Linear probing~\cite{alain2016understanding}, concept activation
vectors~\cite{kim2018interpretability}, network dissection~\cite{bau2017network} and representational
similarity analysis~\cite{kriegeskorte2008representational} were developed for convolutional networks and
language models, and were applied early to medical classifiers~\cite{lucieri2020interpretability}. What is
new is their application to medical foundation models: embedding geometry~\cite{restrepo2026cone},
demographic probing~\cite{zheng2024demographic}, sparse autoencoders for medical encoders and
MLLMs~\cite{bouzid2025insights,nakka2025mammosae,wesp2026sparse}, and steering, editing and patching of
medical multimodal models~\cite{nooralahzadeh2026universal,zhu2026evaluating,sadanandan2026mechanistically}.
Most of this work dates from 2025 onwards (Fig.~\ref{fig:corpus}a) and is spread across machine-learning,
medical-imaging and clinical venues. Existing surveys cover medical explainable AI, medical foundation
models or general interpretability separately; none reviews representation-level evidence across medical
foundation models (Table~\ref{tab:surveys}).

\textbf{Scope.} This survey covers measurements of, and interventions on, the internal representations of
medical vision, vision--language and multimodal foundation models, including their language decoders, and
the general-domain methods they build on. Clinical language models, foundation models for electronic health
records (EHR), physiological signals and omics, and generative medical image models fall outside the core
scope and appear only as marked anchors or instruments. Input attribution enters only when it tests a
representation-level claim (\S\ref{sec:method}). Terminology follows \S\ref{sec:prelim}.

\textbf{Contributions.} The survey makes four contributions. \textit{First}, a taxonomy of feature-space
interpretability methods by family, locus and concept provenance, with rules for the boundary cases
(\S\ref{sec:taxonomy}), applied to \Ncore{} core studies, \Npreprint{} preprints among them, coded at
the level of individual results (\S\ref{sec:method}). \textit{Second}, a family-by-family review of what each family can and cannot
establish, what it costs, and how the families combine (\S\ref{sec:methods}). \textit{Third}, a
consolidated account of datasets, concept resources, benchmarks and protocols, with an assessment of how
comparable the evidence is (\S\ref{sec:evaluation}), and an evidence map of \Fnrows{} statements, each
read against the controls that the studies behind it report (\S\ref{sec:findings}), together with a review
of the experimental uses of these methods (\S\ref{sec:applications}). \textit{Fourth}, the open problems,
and what a solution to each would have to show (\S\ref{sec:future}).

\begin{figure*}[!tbp]
\centering
\includegraphics[width=0.98\linewidth]{fig1.png}
\caption{Overview. The three architectural templates of \S\ref{sec:models} share one computational path:
(A) a unimodal vision encoder feeding a downstream probe or classifier; (B) a dual encoder whose image
and text encoders meet in a projected joint space; (C) a multimodal language model, in which a connector
$\pi$ turns encoder features into visual tokens that join the prompt's text tokens in the decoder's
residual stream and end in the pre-readout state before the logits, answer or report. The labels along
the path are the six loci of \S\ref{sec:loci}. The lower strip names the five method families of
\S\ref{sec:methods}, and the axis beneath them is the evidence-type tag, observational (what is
encoded) against interventional (what is used), which belongs to each claim rather than to a family.}
\label{fig:overview}
\end{figure*}
